\section{AR队列跟随的完整启动命令}\label{app:queue-commands}

将功能包放入\code{~/catkin\_ws/src/}后编译：
\begin{lstlisting}[language=bash]
source /opt/ros/kinetic/setup.bash
cd ~/catkin_ws
catkin_make
source devel/setup.bash
\end{lstlisting}
后车分终端启动底盘、Kinect和深度扫描转换：
\begin{lstlisting}[language=bash]
# TurtleBot2底盘
source /opt/ros/kinetic/setup.bash
source ~/catkin_ws/devel/setup.bash
roslaunch turtlebot_bringup minimal.launch

# Kinect
source /opt/ros/kinetic/setup.bash
source ~/catkin_ws/devel/setup.bash
roslaunch openni_launch openni.launch

# 深度图转换为 /scan
source /opt/ros/kinetic/setup.bash
rosrun depthimage_to_laserscan depthimage_to_laserscan \
  image:=/camera/depth/image_raw \
  camera_info:=/camera/depth/camera_info
\end{lstlisting}
确认\code{/scan}持续更新后，二号和三号车分别使用下列参数启动跟随节点：
\begin{lstlisting}[language=bash]
# 二号车跟随一号车（标签ID 0）
roslaunch team_turtlebot_follow_plus follow_ar.launch \
  marker_id:=0 marker_size:=14.0 desired_distance:=0.40 \
  max_linear_speed:=0.35 max_angular_speed:=0.25 \
  avoid_distance:=0.38 stop_distance:=0.20 side_clearance:=0.40 \
  avoid_release_margin:=0.20 obstacle_turn_speed:=0.20 \
  marker_timeout:=0.8 lost_search_timeout:=16 search_angular_speed:=0.20

# 三号车跟随二号车（标签ID 1）
roslaunch team_turtlebot_follow_plus follow_ar.launch \
  marker_id:=1 marker_size:=14.0 desired_distance:=0.40 \
  max_linear_speed:=0.35 max_angular_speed:=0.25 \
  avoid_distance:=0.38 stop_distance:=0.20 side_clearance:=0.40 \
  avoid_release_margin:=0.20 obstacle_turn_speed:=0.20 \
  marker_timeout:=0.8 lost_search_timeout:=16 search_angular_speed:=0.20
\end{lstlisting}
